\section*{3. 硬币找零问题}

假设有数目不限的面值为25美分、10美分、5美分、1美分的硬币，请使用最少个数的硬币凑出3.33美元。

\begin{itemize}
    \item \textbf{贪心策略}：优先使用面值\textbf{最大}的硬币
    \item \textbf{注意事项}：此策略对特定面值体系（如美分）有效，对任意面值未必有效（需用动态规划）
    \item \textbf{时间复杂度}：$O(K)$，K为硬币种类数
\end{itemize}

\begin{lstlisting}[language=C++, style=cppstyle]
void solve_coins(int cents) {
    int coins[] = {25, 10, 5, 1};
    int count = 0;

    for(int val : coins) {
        int num = cents / val;
        count += num;
        cents %= val;
        cout << "面值 " << val << ": "
             << num << " 个" << endl;
    }
    cout << "总个数: " << count << endl;
}
\end{lstlisting}

\textbf{解题步骤：}
\begin{enumerate}
    \item 将硬币面值按从大到小排序
    \item 从最大面值开始，计算当前金额能包含多少个该面值硬币
    \item 用取模运算得到剩余金额
    \item 对下一个面值重复上述过程，直到金额为0
    \item 累加所有硬币数量得到答案
\end{enumerate}

\textbf{示例：} 333美分 $= 13 \times 25 + 0 \times 10 + 1 \times 5 + 3 \times 1 = 17$个硬币
